% Generated draft: requires mathematical and source review.
\documentclass[11pt]{article}
\usepackage[margin=25mm]{geometry}
\usepackage{fontspec}
\setmainfont{Latin Modern Roman}
\usepackage{amsmath,amssymb,mathtools,unicode-math}
\setmathfont{Latin Modern Math}
\usepackage{enumitem,needspace,xurl,hyperref}
\setlength{\parindent}{0pt}
\setlength{\parskip}{5pt}
\setlength{\emergencystretch}{3em}
\newcommand{\N}{\mathbb{N}}
\newcommand{\Z}{\mathbb{Z}}
\newcommand{\Q}{\mathbb{Q}}
\newcommand{\R}{\mathbb{R}}
\newcommand{\C}{\mathbb{C}}
\newcommand{\F}{\mathbb{F}}
\newcommand{\A}{\mathbb{A}}
\newcommand{\PP}{\mathbb{P}}
\newcommand{\E}{\mathbb{E}}
\newcommand{\eps}{\varepsilon}
\newcommand{\dd}{\,\mathrm{d}}
\begin{document}
\section*{UnsolvedMath ID: 701514}
\textbf{Source problem code:} catalogue entry
\subsection{Definitions and mathematical statement}
Let \(G\) be a group with identity \(1_G\). For subsets \(A,B\subseteq G\), define
\[
AB=\{ab:a\in A,\ b\in B\},\qquad A^*=\bigcup_{n\geq 0}A^n,
\]
where \(A^0=\{1_G\}\) and \(A^{n+1}=A^nA\). Thus \(A^*\) is the submonoid generated by \(A\). The family \(\operatorname{Rat}(G)\) of rational subsets of \(G\) is the smallest family containing every finite subset of \(G\) and closed under finite unions, setwise products, and the operation \(A\mapsto A^*\). To say that this family forms a Boolean algebra means that it is closed under finite unions, finite intersections, and complements relative to \(G\).

A group is finitely generated if some finite subset generates it as a group. Define its derived series by \(G^{(0)}=G\) and \(G^{(n+1)}=[G^{(n)},G^{(n)}]\), where \([H,H]\) is the subgroup generated by the commutators \(x^{-1}y^{-1}xy\) with \(x,y\in H\). The group is solvable if \(G^{(d)}=\{1_G\}\) for some integer \(d\geq 0\). It is virtually abelian if it has an abelian subgroup of finite index.

\textbf{Conjecture (V. A. Roman'kov; Kourovka Notebook 18.85).} For every finitely generated solvable group \(G\),
\[
\operatorname{Rat}(G)\text{ forms a Boolean algebra}
\quad\Longrightarrow\quad
\exists H\leq G\;\bigl([G:H]<\infty\ \text{and}\ H\text{ is abelian}\bigr).
\]

The supplied source reports that the implication is known when \(G\) is metabelian, polycyclic, or of finite rank, and attributes these special cases to G. A. Bazhenova. Here metabelian means \(G^{(2)}=\{1_G\}\), and polycyclic means that \(G\) has a finite subnormal series with cyclic factors. The source does not specify its convention for finite rank; that qualification is retained without choosing a definition. These special-case claims and the record's open status are supplied source assessments, not independently verified conclusions.
\subsection{Short English statement}
Roman'kov's conjecture asks whether every finitely generated solvable group whose rational subsets form a Boolean algebra has an abelian subgroup of finite index. The supplied record labels the problem open and reports special cases attributed to G. A. Bazhenova; its meaning of finite rank is unspecified.
\subsection{Sources}
\begin{itemize}
\item Supplied local UnsolvedMath registry record: permanent ID \(701514\), exact source problem code \texttt{catalogue entry}, title \emph{Kourovka Notebook Problem 18.85}. The record labels the problem open and records V. A. Roman'kov as proposer, with proposed year 2014.
\item \emph{The Kourovka Notebook (Unsolved Problems in Group Theory)}, original problem 18.85, 18th issue (2014); the supplied background also identifies the 21st issue edition (2026), \texttt{arXiv:1401.0300v48}. The supplied collection description names E. I. Khukhro and V. D. Mazurov as editors. This bibliographic attribution is supplied evidence and has not been independently checked.
\item The supplied statement attributes the metabelian, polycyclic, and finite-rank special cases to G. A. Bazhenova, but supplies no separate publication citation for them.
\end{itemize}
\end{document}
